\exercisesection{2}

\begin{enumerate}
\def\labelenumi{\arabic{enumi}.}
\tightlist
\item
  A multiplicative subset \(S\) of a ring \(A\) is called saturated if the relation \(xy \in S\) implies \(x \in S\) and \(y \in S\) .
\end{enumerate}

\begin{enumerate}
\def\labelenumi{(\alph{enumi})}
\item
  For a subset \(S\) of \(A\) to be multiplicative and saturated, it is necessary and sufficient that \(A - S\) be a union of prime ideals of \(A\) .
\item
  Let S be a multiplicative subset of A and \(\tilde{S}\) the set of \(x \in A\) for which there exists \(y \in A\) such that \(xy \in S\) . Show that \(\tilde{S}\) is the smallest saturated multiplicative subset containing S, A --- \(\tilde{S}\) is the union of the prime ideals of A not meeting S and, for every A-module M, the canonical mapping from \(S^{-1}M\) to \(\tilde{S}^{-1}M\) is bijective.
\item
  Let S and T be two multiplicative subsets of A. Show that the two following properties are equivalent: (α) \(S \subset \tilde{T}\) ; (β) for every A-module M such that \(S^{-1}M = 0\) , \(T^{-1}M = 0\) . (To see that (β) implies (α), consider a quotient A-module A/As, where \(s \in S\) .)
\end{enumerate}

\begin{enumerate}
\def\labelenumi{\arabic{enumi}.}
\setcounter{enumi}{1}
\tightlist
\item
  Let A be a ring and S a multiplicative subset of A. Show that the set n of \(x \in A\) such that sx = 0 for some \(s \in S\) is an ideal of A. Set \(A_{1} = A/n\) and denote by \(S_{1}\) the canonical image of S in \(A_{1}\) . Show that no element of \(S_{1}\) is a divisor of 0 in \(\mathbf{A}_1\) and that the canonical homomorphism \(\mathbf{S}^{-1}\mathbf{A} \to \mathbf{S}_1^{-1}\mathbf{A}_1\) is bijective.
\end{enumerate}

Deduce that, for \(S^{-1}A\) to be a finitely generated A-module, it is necessary and sufficient that all the elements of \(S_{1}\) be invertible in \(A_{1}\) , in which case \(S^{-1}A\) is identified with \(A_{1}\) .

\begin{enumerate}
\def\labelenumi{\arabic{enumi}.}
\setcounter{enumi}{2}
\tightlist
\item
  Let \(S = Z^{*}\) (the complement of \(\{0\}\) in Z), p an integer \textgreater1 and M the Z-module the direct sum of the modules \(Z/p^{n}Z\) for \(n \in N\) . Show that \(S^{-1}M = 0\) although M is a faithful Z-module and that the canonical mapping
\end{enumerate}

\[
\mathrm{S} ^ {- 1} \operatorname{End} _ {\mathbf {z}} (\mathrm{M}) \rightarrow \operatorname{End} _ {\mathrm{s} ^ {- 1} \mathbf {z}} (\mathrm{S} ^ {- 1} \mathrm{M})
\]

is not injective.

\begin{enumerate}
\def\labelenumi{\arabic{enumi}.}
\setcounter{enumi}{3}
\item
  Let \(S = \mathbf{Z}^*\) and \(M\) be a free \(\mathbf{Z}\) -module with infinite basis. Show that the canonical mapping \(S^{-1}\operatorname{End}_{\mathbf{Z}}(M) \to \operatorname{End}_{S^{-1}\mathbf{Z}}(S^{-1}M)\) is not surjective. Deduce from this and Exercise 3 an example of a \(\mathbf{Z}\) -module \(M\) such that the canonical mapping \(S^{-1}\operatorname{End}_{\mathbf{Z}}(M) \to \operatorname{End}_{S^{-1}\mathbf{Z}}(S^{-1}M)\) is neither injective nor surjective.
\item
  Give an example of a sequence \((\mathbf{P}_n)\) of sub- \(\mathbf{Z}\) -modules of \(\mathbf{Z}\) such that, for \(\mathbf{S} = \mathbf{Z}^*\) , the submodule \(\mathrm{S}^{-1}\left(\bigcap_{n} \mathrm{P}_{n}\right)\) is distinct from \(\bigcap_{n} (\mathrm{S}^{-1}\mathrm{P}_{n})\) .
\item
  Give an example of a \(\mathbf{Z}\) -module \(\mathbf{M}\) such that, for \(\mathbf{S} = \mathbf{Z}^*\) , \(\operatorname{Ann}(\mathbf{M}) = 0\) , but \(\operatorname{Ann}(\mathbf{S}^{-1}\mathbf{M}) = \mathbf{Q}\) (cf.~Exercise 3).
\item
  Let S be a multiplicative subset of a ring A; for every family \((M_{t})_{t\in I}\) of A-modules, define a canonical homomorphism \(S^{-1}\prod_{i\in I}M_{i}\to\prod_{i\in I}S^{-1}M_{i}\) of \(S^{-1}A\) -modules. Give an example where this homomorphism is neither injective nor surjective (cf.~Exercise 3).
\item
  Let A be a ring, \((\mathrm{S}_{\lambda})_{\lambda \in L}\) a family of multiplicative subsets of A and \(M = \bigoplus_{\lambda \in L} S_{\lambda}^{-1}A\) . Show that the two following properties are equivalent: ( \(\alpha\) ) M is a faithfully flat A-module; ( \(\beta\) ) for every maximal ideal m of A, there exists \(\lambda \in L\) such that \(m \cap S_{\lambda} = \varnothing\) . In particular, if S is a multiplicative subset of A, the A-module \(S^{-1}A\) is faithfully flat only if the elements of S are invertible, in which case \(S^{-1}A\) is identified with A.
\item
  Let A be a ring, S a multiplicative subset of A and q a prime ideal of A. Show that, if T is the image of A - q in \(S^{-1}A\) , the ring \(T^{-1}(S^{-1}A)\) is isomorphic to \(S'^{-1}A\) , where \(S'\) is the complement of the union of the prime ideals of A contained in q and not meeting S.
\item
  \begin{enumerate}
  \def\labelenumii{(\alph{enumii})}
  \tightlist
  \item
    Let K be a commutative field, A the polynomial ring K{[}X, Y{]}, \(p = (X)\) , \(q = (Y)\) and S the complement of \(p \cup q\) in A. Show that the saturation of \(p + q\) with respect to S is distinct from the sum of the saturations of p and q.
  \end{enumerate}
\end{enumerate}

\begin{enumerate}
\def\labelenumi{(\alph{enumi})}
\setcounter{enumi}{1}
\tightlist
\item
  Let K be a commutative field, A the polynomial ring K{[}X, Y, Z{]}, \(\mathfrak{p} = (\mathbf{X}) + (\mathbf{Z})\) , \(\mathfrak{q} = (\mathbf{Y}) + (\mathbf{Z})\) and S the complement of \(p \cup q\) in A. Show that the saturation of pq with respect to S is distinct from the product of the saturations of p and q.
\end{enumerate}

\begin{enumerate}
\def\labelenumi{\arabic{enumi}.}
\setcounter{enumi}{10}
\item
  Let \(p\) be a prime number; which ideals of \(\mathbf{Z}\) are saturated with respect to the ideal \((p)\) ?
\item
  In a ring A let \(\mathfrak{r}(\mathfrak{a})\) denote the radical of an ideal \(\mathfrak{a}\) .
\end{enumerate}

\begin{enumerate}
\def\labelenumi{(\alph{enumi})}
\tightlist
\item
  Show that, for three ideals a, b, c of A,
\end{enumerate}

\[
\mathfrak {r} (\mathfrak {a} + \mathfrak {b c}) = \mathfrak {r} (\mathfrak {a} + (\mathfrak {b} \cap \mathfrak {c})) = \mathfrak {r} (\mathfrak {a} + \mathfrak {b}) \cap \mathfrak {r} (\mathfrak {a} + \mathfrak {c}).
\]

\begin{enumerate}
\def\labelenumi{(\alph{enumi})}
\setcounter{enumi}{1}
\tightlist
\item
  Give an example of an infinite sequence of ideals \((\mathfrak{a}_n)\) of A such that \(\mathfrak{r}\left(\bigcap_{n} \mathfrak{a}_n\right) \neq \bigcap_{n} \mathfrak{r}(\mathfrak{a}_n)\) .
\end{enumerate}

*(c) Give an example of a principal ideal \(\mathfrak{a}\) such that \(\mathfrak{r}(\mathfrak{a})\) is not finitely generated and there exists no integer \(n > 0\) such that \((\mathfrak{r}(\mathfrak{a}))^n \subset \mathfrak{a}\) . (Consider the ring of a non-discrete valuation of height 1.)*

13. Let A be a ring.\par

\begin{enumerate}
\def\labelenumi{(\alph{enumi})}
\item
  Give another proof of Exercise 6 (b) of Algebra, Chapter VIII, § 6, by considering the image of the polynomial \(f\) in the rings \((\mathbf{A} / \mathfrak{p})[\mathbf{X}]\) , where \(\mathfrak{p}\) runs through the set of prime ideals of \(\mathbf{A}\) .
\item
  Let \(N\) be a square matrix of order \(r\) over \(A\) ; for \(N\) to be nilpotent, it is necessary and sufficient that each of the coefficients (other than the dominant coefficient) of its characteristic polynomial be so. (To see that the condition is necessary, use the same method as in (a); to see that it is sufficient, use the Cayley-Hamilton Theorem.)
\item
  Show that, for every ordered pair of positive integers \(k, r\) , there exists a least positive integer \(m(k, r)\) independent of the ring \(\mathbf{A}\) such that the relation \(N^k = 0\) implies \((\operatorname{Tr}(N))^{m(k,r)} = 0\) . (Consider the ring \(\Lambda = \mathbf{Z}[T_{ij}]\) , where the \(T_{ij}\) are \(r^2\) indeterminates, the matrix \(X = (T_{ij})\) over \(\Lambda\) and the element \(t = \sum_{i=1}^{n} \mathrm{T}_{ii} = \mathrm{Tr}(X)\) of \(\Lambda\) ; if \(\mathfrak{a}_k\) is the ideal of \(\Lambda\) generated by the elements of \(X^k\) , show, with the aid of (b), that there exists an integer \(m\) such that \(t^m \in \mathfrak{a}_k\) .)
\item
  Show that \(m(2, 2) = 4\) .
\end{enumerate}

¶ 14. (a) Let A be a left Noetherian ring (not necessarily commutative) and a a left ideal of A. Suppose that there exists a left ideal b ≠ 0 such that a ∩ b = 0.

Show that every element \(a \in \mathfrak{a}\) is a right divisor of 0. (Let \(b \neq 0\) be an element of b and let \(a_{n} = Ab + Aba + \cdots + Aba^{n}\) ; consider the smallest integer n such that \(a_{n} = a_{n+1}\) .

\begin{enumerate}
\def\labelenumi{(\alph{enumi})}
\setcounter{enumi}{1}
\tightlist
\item
  Let E be an infinite-dimensional vector space over a field K and A the ring \(\operatorname{End}_{\mathbb{K}}(\operatorname{E})\) . Give an example of two elements u, v of A such that neither u nor v is a right divisor of 0 and \(Au \cap Av = 0\) .
\end{enumerate}

¶ 15. Let X be a completely regular topological space and βX its Stone-Čech compactification (General Topology, Chapter IX, § 1, Exercise 7). For all \(x \in \mathbf{X}\) , let \(\mathfrak{J}(x)\) denote the maximal ideal of the ring \(\mathcal{C}(\mathbf{X}; \mathbf{R})\) consisting of those \(f \in \mathcal{C}(\mathbf{X}; \mathbf{R})\) such that \(x\) belongs to the closure of the sub set \(\bar{f}^1(0)\) of X (General Topology, Chapter X, § 4, Exercise 15); let \(\mathfrak{U}(x)\) denote the ideal consisting of those \(f \in \mathcal{C}(\mathbf{X}; \mathbf{R})\) such that \(\bar{f}^1(0)\) is the trace on X of a neighbourhood of \(x\) in X; \(\mathfrak{U}(x)\) is equal to its radical.

An ideal \(\mathfrak{a}\) of \(\mathfrak{C}(\mathbf{X};\mathbf{R})\) is called isolated (resp. absolutely isolated) if, for all \(f\geqslant 0\) on \(\mathfrak{a}\) , every \(g\) such that \(0\leqslant g\leqslant f\) (resp. \(|g|\leqslant f\) ) belongs to \(\mathfrak{a}\) (cf.~Algebra, Chapter IV, §1, Exercise 4); then \(\mathcal{C}(\mathbf{X};\mathbf{R}) / \mathfrak{a}\) possesses an order structure (resp. a lattice-order structure) compatible with its ring structure and for which the elements \(\geqslant 0\) are the classes mod. \(\mathfrak{a}\) of the elements \(\geqslant 0\) to \(\mathcal{C}(\mathbf{X};\mathbf{R})\) (loc. cit.).

\begin{enumerate}
\def\labelenumi{(\alph{enumi})}
\item
  Show that an ideal \(\mathfrak{a}\) of \(\mathcal{C}(\mathbf{X};\mathbf{R})\) which is equal to its radical is absolutely isolated (if \(f\in \mathfrak{a}\) and \(|g|\leqslant |f|\) , consider the function equal to 0 when \(f(x) = 0\) and to \((g(x))^2 /f(x)\) when \(f(x)\neq 0\) ). In this case for \(\mathcal{C}(\mathbf{X};\mathbf{R}) / \mathfrak{a}\) to be totally ordered, it is necessary and sufficient that \(\mathfrak{a}\) be prime; the set of prime ideals of \(\mathcal{C}(\mathbf{X};\mathbf{R}) / \mathfrak{a}\) is then totally ordered by inclusion.
\item
  For an isolated ideal \(\mathfrak{a}\) of \(\mathcal{C}(\dot{\mathbf{X}};\mathbf{R})\) to be such that \(\mathcal{C}(\mathbf{X};\mathbf{R}) / \mathfrak{a}\) is totally ordered, it is necessary that \(\mathfrak{a}\) be contained in a single maximal ideal \(\mathfrak{J}(x)\) and then \(\mathfrak{a}\) necessarily contains \(\mathfrak{U}(x)\) ; it is sufficient that \(\mathfrak{a}\) contain a prime ideal and then \(\mathfrak{a}\) is absolutely isolated (note that the relation \(fg = 0\) then implies \(f\in \mathfrak{a}\) or \(g\in \mathfrak{a}\) ).
\item
  Take \(\mathbf{X} = \mathbf{R}\) , \(x = 0\) ; show that the principal ideal generated by the function \(t \mapsto |t|\) is contained in \(\mathfrak{J}(0)\) and contains \(\mathfrak{U}(0)\) but is not isolated; the principal ideal generated by the identity mapping \(t \mapsto t\) is contained in \(\mathfrak{J}(0)\) , contains \(\mathfrak{U}(0)\) and is isolated but not absolutely isolated.
\item
  Show that, if \(\mathfrak{p}\) and \(\mathfrak{q}\) are two prime ideals in \(\mathcal{C}(\mathbf{X};\mathbf{R})\) , then \(\mathfrak{pq} = \mathfrak{p} \cap \mathfrak{q}\) (and in particular \(\mathfrak{p}^2 = \mathfrak{p}\) ) and \(\mathfrak{p} + \mathfrak{q}\) is prime or equal to \(\mathcal{C}(\mathbf{X};\mathbf{R})\) . (If \(\mathfrak{p} + \mathfrak{q} \neq \mathcal{C}(\mathbf{X};\mathbf{R})\) , observe that \(\mathfrak{p} + \mathfrak{q}\) is absolutely isolated and equal to its radical and use (a).)
\item
  For \(\mathfrak{U}(x)\) to be prime, it is necessary and sufficient that, for every function \(f\in\mathcal{C}(\mathbf{X};\mathbf{R})\) , the closure in \(\beta\mathbf{X}\) of one of the two sets \(\bar{f}^{-1}\left([0,+\infty[),\bar{f}^{-1}([- \infty,0])\right)\) is a neighbourhood of \(x\) . Deduce an example where \(\mathfrak{U}(x)\) is prime and not maximal (consider the topological space associated with a non-trivial ultrafilter; cf.~General Topology, Chapter I, § 6, no. 5).
\item
  Suppose that \(x \in \mathbf{X}\) . For \(\mathfrak{U}(x) = \mathfrak{J}(x)\) , it is necessary and sufficient that every countable intersection of neighbourhoods of \(x\) in \(\mathbf{X}\) be a neighbourhood of \(x\) in \(\mathbf{X}\) . (Cf. General Topology, Chapter I, § 7, Exercise 7.)
\item
  Show that \(\mathcal{C}(\mathbf{X};\mathbf{R})\) is canonically identified with the total ring of fractions of \(\mathcal{C}^{\infty}(\mathbf{X};\mathbf{R})\) .
\end{enumerate}

\begin{enumerate}
\def\labelenumi{\arabic{enumi}.}
\setcounter{enumi}{15}
\tightlist
\item
  Let A be a topological ring (not necessarily commutative). The topology on A is called (left) linear if the open left ideals of A form a fundamental system of neighbourhoods of 0. If this is so, the set F of open left ideals of A satisfies the following conditions:
\end{enumerate}

\begin{enumerate}
\def\labelenumi{(\arabic{enumi})}
\item
  Every left ideal of A containing an ideal \(m \in \mathcal{F}\) belongs to \(\mathcal{F}\) .
\item
  Every finite intersection of ideals of \(\mathcal{F}\) belongs to \(\mathcal{F}\) .
\item
  If \(\mathfrak{m} \in \mathcal{F}\) and \(a \in \mathbf{A}\) , the left ideal \(\mathfrak{n}\) consists of those \(x \in \mathbf{A}\) such that \(xa \in \mathfrak{m}\) belongs to \(\mathcal{F}\) .
\end{enumerate}

Conversely, if a set F of left ideals of a ring A satisfies these three conditions, there exists a unique topology compatible with the ring structure on A and for which F is the set of open left ideals of A. Such a set F of left ideals is called (left) topologizing.

\begin{enumerate}
\def\labelenumi{\arabic{enumi}.}
\setcounter{enumi}{16}
\tightlist
\item
  (*) Let A be a ring (not necessarily commutative) and \(\mathcal{F}\) a topologizing set (Exercise 16) of left ideals of A.
\end{enumerate}

\begin{enumerate}
\def\labelenumi{(\alph{enumi})}
\item
  A (left) A-module M is called F-negligible if the annihilator of every element of M belongs to F. If M is F-negligible, so is every submodule and every quotient module of M. If M and N are two F negligible A-modules, \(M \oplus N\) is F-negligible. For \(A_{s}/m\) to be F-negligible, it is necessary and sufficient that \(m \in F\) . If \((0) \in F\) , every A-module is F-negligible; if consists only of the single left ideal \(A_{s}\) , every F-negligible A-module is reduced to 0.
\item
  For every A-module M there exists a greatest F-negligible submodule of M which is denoted by FM. For every A-module homomorphism \(u: M \to N\) , \(u(\mathcal{F}M) \subset \mathcal{F}N\) ; let Fu be the mapping from FM to FN whose graph coincides with that of \(u|FM\) . If
\end{enumerate}

\[
0 \longrightarrow \mathrm{M} \stackrel {u} {\longrightarrow} \mathrm{N} \stackrel {v} {\longrightarrow} \mathrm{P}
\]

is an exact sequence of A-modules, the sequence

\[
0 \longrightarrow \mathcal {F} M \stackrel {\mathcal {F} _ {u}} {\longrightarrow} \mathcal {F} N \stackrel {\mathcal {F} _ {v}} {\longrightarrow} \mathcal {F} P
\]

is exact. Show that \(FA_{s}\) is a two-sided ideal of A.

\begin{enumerate}
\def\labelenumi{(\alph{enumi})}
\setcounter{enumi}{2}
\tightlist
\item
  Let \(\mathfrak{m}\) , \(\mathfrak{n}\) be two elements of \(\mathcal{F}\) such that \(\mathfrak{m} \subset \mathfrak{n}\) ; for every left A-module \(\mathbf{M}\) the canonical injection \(j \colon \mathfrak{m} \to \mathfrak{n}\) defines a commutative group homomorphism
\end{enumerate}

\[
u _ {m, n} = \operatorname{Hom} _ {A} (j, 1 _ {M}): \operatorname{Hom} _ {A} (n, M) \rightarrow \operatorname{Hom} _ {A} (m, M).
\]

If \(\mathcal{F}\) is ordered by the relation \(\supset\) , the \(u_{m,n}\) define a direct system of commutative groups, whose direct limit will be denoted by \(\mathrm{M}_{(\mathcal{F})}\) . The mappings \(u_{m,A_s}\) form a direct system, whose direct limit is a mapping (called canonical) from \(\mathrm{M} = \mathrm{Hom}_{\mathrm{A}}(\mathrm{A}_s,\mathrm{M})\) to \(\mathrm{M}_{(\mathcal{F})}\) . Similarly, if \(v: \mathrm{M} \to \mathrm{N}\) is a left A-module homomorphism, the \(v_m = \mathrm{Hom}_{\mathrm{A}}(1,v): \mathrm{Hom}_{\mathrm{A}}(\mathfrak{m},\mathrm{M}) \to \mathrm{Hom}_{\mathrm{A}}(\mathfrak{m},\mathrm{N})\) form a direct system, whose direct limit is a commutative group homomorphism

\[
v _ {(\mathcal {F})} \colon \mathrm{M} _ {(\mathcal {F})} \rightarrow \mathrm{N} _ {(\mathcal {F})}.
\]

If \(0 \to \mathbf{M} \xrightarrow{v} \mathbf{N} \xrightarrow{w} \mathbf{P}\) is an exact sequence of A-modules, the sequence

\[
0 \longrightarrow \mathrm{M} _ {(\mathcal {F})} \stackrel {v (\mathcal {F})} {\longrightarrow} \mathrm{N} _ {(\mathcal {F})} \stackrel {w (\mathcal {F})} {\longrightarrow} \mathrm{P} _ {(\mathcal {F})}
\]

is exact.

\begin{enumerate}
\def\labelenumi{\arabic{enumi}.}
\setcounter{enumi}{17}
\tightlist
\item
  \begin{enumerate}
  \def\labelenumii{(\alph{enumii})}
  \tightlist
  \item
    Let A be a ring and F and G two topologizing sets (Exercise 16) of left ideals of A. Let F.G denote the set of left ideals l of A for which there exists a left ideal n ∈ G such that n/l is F-negligible. Show that F.G is topologizing; for an A-module M to be F.G-negligible, it is necessary and sufficient that there exist an F-negligible submodule M' of M such that M/M' is G-negligible. Show that the composition law (F,G) → F.G on the set of topologizing sets of left ideals of A is associative.
  \end{enumerate}
\end{enumerate}

\(\mathcal{F}\) is called idempotent if \(\mathcal{F}.\mathcal{F} = \mathcal{F}\) .

\begin{enumerate}
\def\labelenumi{(\alph{enumi})}
\setcounter{enumi}{1}
\item
  Show that \(\mathcal{F}.\mathcal{G}\) contains all the ideals of the form \(\mathfrak{m}.\mathfrak{n}\) where \(\mathfrak{m}\in\mathcal{F}\) and \(\mathfrak{n}\in\mathcal{G}\) .
\item
  Let K be a commutative field and B the ring of polynomials with coefficients in K in an infinite system of indeterminates \((\mathbf{X}_{i})\) ( \(i \geqslant 0\) ). Let b be the ideal of B generated by the elements \(X_{i}X_{j}\) ( \(i \neq j\) ); let A be the ring B/b, \(\xi_{i}\) the class mod. b of \(X_{i}\) and m the ideal of A generated by the \(\xi_{i}\) ; take F to be the set of ideals containing a power of m. Show that F is topologizing and contains the product of any two ideals of F but is not idempotent.
\item
  Suppose that A is commutative and F is a topologizing set of ideals of A such that every ideal \(m \in F\) contains a finitely generated ideal \(n \in F\) . Show that, if the product of two ideals of F belongs to F, F is idempotent.
\end{enumerate}

¶ 19. Let F be a set of left ideals of a ring A; suppose that F is topologizing and idempotent (Exercises 16 and 18).

\begin{enumerate}
\def\labelenumi{(\alph{enumi})}
\tightlist
\item
  Show that, if \(\mathfrak{m} \in \mathcal{F}\) , \(\mathfrak{n} \in \mathcal{F}\) and \(u \in \operatorname{Hom}_{\mathbf{A}}(\mathfrak{n}, \mathbf{A}_s)\) , then \(\bar{u}^1(\mathfrak{m}) \in \mathcal{F}\) (consider the exact sequence \(0 \to \mathfrak{n} / \bar{u}^1(\mathfrak{m}) \to \mathbf{A}_s / \bar{u}^1(\mathfrak{m}) \to \mathbf{A}_s / \mathfrak{n} \to 0\) ). For every A-module M and all \(v \in \operatorname{Hom}_{\mathbf{A}}(\mathfrak{m}, \mathbf{M})\) , let \(u.v\) denote the canonical image in \(\mathbf{M}_{(\mathcal{F})}\) (Exercise 17 (c)) of the composite homomorphism \(\bar{u}^1(\mathfrak{m}) \xrightarrow{u} \mathfrak{m} \xrightarrow{v} \mathbf{M}\) . Show that the mapping \((u,v) \mapsto u.v\) is Z-bilinear; there corresponds to it a Z-linear mapping
\end{enumerate}

\[
\psi_ {m, n} \colon \operatorname{Hom} _ {A} (n, A _ {s}) \otimes_ {\mathbf {Z}} \operatorname{Hom} _ {A} (m, M) \rightarrow M _ {(\mathcal {F})}.
\]

These mappings form a direct system and their direct limit therefore corresponds canonically to a Z-bilinear mapping

\[
\psi \colon A _ {(\mathcal {F})} \times M _ {(\mathcal {F})} \rightarrow M _ {(\mathcal {F})}.
\]

Show that, if \(M = A_{s}\) , \(\psi\) is an internal law of composition which makes \(\mathbf{A}_{(\mathcal{F})}\) into a ring and the canonical mapping \(A \to A_{(\mathcal{F})}\) is a ring homomorphism for this structure. For any left A-module M, \(\psi\) is an external law defining on \(\mathbf{M}_{(\mathcal{F})}\) a left \(\mathbf{A}_{(\mathcal{F})}\) -module structure; with this structure (and the canonical mapping \(A \to A_{(\mathcal{F})}\) ) the canonical mapping \(M \to M_{(\mathcal{F})}\) is an A-module homomorphism.

\begin{enumerate}
\def\labelenumi{(\alph{enumi})}
\setcounter{enumi}{1}
\tightlist
\item
  If \(i: \mathbf{M} \to \mathbf{M}_{(\mathcal{F})}\) is the canonical homomorphism, show that
\end{enumerate}

\[
\operatorname{Ker} (i) = \mathcal {F M}
\]

and that the A-module Coker \((i) = \mathbf{M}_{(\mathcal{F})} / i(\mathbf{M})\) is \(\mathcal{F}\) -negligible.

\begin{enumerate}
\def\labelenumi{(\alph{enumi})}
\setcounter{enumi}{2}
\tightlist
\item
  For every left A-module M let \(M_{F}\) denote the left A-module \((\mathrm{M}/\mathcal{F}\mathrm{M})_{(\mathcal{F})}\) and \(j_{M}\) the composite canonical mapping
\end{enumerate}

\[
\mathbf {M} \rightarrow \mathbf {M} / \mathcal {F} \mathbf {M} \rightarrow \mathbf {M} _ {\mathcal {F}};
\]

then \(\operatorname{Ker}(j_{\mathbf{M}}) = \mathcal{F}\mathbf{M}\) and the A-module \(\operatorname{Coker}(j_{\mathbf{M}}) = \mathrm{M}_{\mathcal{F}} / j_{\mathbf{M}}(\mathbf{M})\) is \(\mathcal{F}\) -negligible.

Let \(u: \mathrm{P} \to \mathrm{M}\) , \(h: \mathrm{P} \to \mathrm{Q}\) be A-module homomorphisms such that \(\operatorname{Ker}(h)\) and \(\operatorname{Coker}(h)\) are \(\mathcal{F}\) -negligible; show that there exists a unique A-homomorphism \(v: \mathrm{Q} \to \mathrm{M}_{\mathcal{F}}\) which makes the following diagram commutative

\[
\begin{array}{c} \mathrm{P} \xrightarrow {u} \mathrm{M} \\ h \Big \downarrow \qquad \qquad \qquad \qquad \qquad \qquad \qquad \Big \downarrow j _ {\mathrm{M}} \\ \mathrm{Q} \xrightarrow [ v ]{} \mathrm{M} _ {\mathcal {F}} \end{array}
\]

(First reduce it to the case where h is injective; then, identifying P with a submodule of Q, note that, for \(x \in Q\) , there exists \(m \in F\) such that \(mx \in P\) and map x to the canonical image in \(M_{F}\) of the composite homomorphism

\[
\mathfrak {m} \xrightarrow {x} \mathfrak {m} x \to \mathrm{P} \xrightarrow {u} \mathrm{M} \to \mathrm{M} / \mathcal {F} \mathrm{M.})
\]

\begin{enumerate}
\def\labelenumi{(\alph{enumi})}
\setcounter{enumi}{3}
\tightlist
\item
  For every A-module homomorphism \(u: M \rightarrow N\) show that there exists a unique homomorphism \(u_{F}: M_{F} \rightarrow N_{F}\) which makes the following diagram commutative
\end{enumerate}

\[
\begin{array}{c c c} \mathrm{M} & \stackrel {{u}} {{\longrightarrow}} \mathrm{N} \\ j _ {\mathrm{M}} \Big \downarrow & & \Big \downarrow j _ {\mathrm{N}} \\ \mathrm{M} _ {\mathcal {F}} & \stackrel {{u _ {\mathcal {F}}}} {{\longrightarrow}} \mathrm{N} _ {\mathcal {F}} \end{array}
\]

(use (a), (b) and (c)). Show that, if

\[
0 \longrightarrow \mathrm{M} \stackrel {u} {\longrightarrow} \mathrm{N} \stackrel {v} {\longrightarrow} \mathrm{P}
\]

is an exact sequence of A-modules, the sequence

\[
0 \longrightarrow \mathrm{M} _ {\mathcal {F}} \stackrel {u \mathcal {F}} {\longrightarrow} \mathrm{N} _ {\mathcal {F}} \stackrel {v \mathcal {F}} {\longrightarrow} \mathrm{P} _ {\mathcal {F}}
\]

is exact (note that the mapping \(\mathbf{M} / \mathcal{F}\mathbf{M}\to \mathrm{N} / \mathcal{F}\mathrm{N}\) derived from \(u\) by taking quotients is injective).

Show that the mappings \(j_{\mathbf{M}_{\mathcal{F}}}\) and \((j_{\mathbf{M}})_{\mathcal{F}}\) coincide and are isomorphisms (and therefore \(\mathcal{FM}_{\mathcal{F}} = 0\) ). If \(\mathbf{N}'\) is a sub-A-module of \(\mathbf{M}_{\mathcal{F}}\) , \((\overline{j_{\mathbf{M}}^{-1}(\mathbf{N}')})_{\mathcal{F}}\) is canonically identified with \(\mathbf{N}_{\mathcal{F}}'\) (observe that, if we write \(\mathbf{N} = j_{\mathbf{M}}^{-1}(\mathbf{N}')\) , \(\mathbf{N}' / j_{\mathbf{M}}(\mathbf{N})\) is \(\mathcal{F}\) -negligible).

\begin{enumerate}
\def\labelenumi{(\alph{enumi})}
\setcounter{enumi}{4}
\tightlist
\item
  Let M be a left A-module and \(\phi: \mathbf{A} \times \mathbf{M} \to \mathbf{M}\) the Z-bilinear mapping \((a, m) \mapsto am\) . Show that there exists a unique Z-bilinear mapping \(\phi_{\mathcal{F}}: \mathbf{A}_{\mathcal{F}} \times \mathbf{M}_{\mathcal{F}} \to \mathbf{M}_{\mathcal{F}}\) which makes the following diagram commutative
\end{enumerate}

\[
\begin{array}{c} \mathbf {A} \times \mathbf {M} \xrightarrow {\phi} \mathbf {M} \\ j _ {\mathbf {A}} \times j _ {\mathbf {M}} \Big \downarrow \qquad \qquad \qquad \qquad \Big \downarrow j _ {\mathbf {M}} \\ \mathbf {A} _ {\mathcal {F}} \times \mathbf {M} _ {\mathcal {F}} \xrightarrow [ \phi_ {\mathcal {F}} ]{} \mathbf {M} _ {\mathcal {F}} \end{array}
\]

We also define on \(A_{F}\) a ring structure and on \(M_{F}\) a left \(A_{F}\) -module structure. If \(u: M \to N\) is an A-module homomorphism, \(u_{F}: M_{F} \to N_{F}\) is an \(A_{F}\) -module homomorphism. In order that, for every A-module M, the mapping \(u \mapsto u_{F}\) of \(\operatorname{Hom}_{\mathbf{A}}(\mathbf{M}, \mathbf{N})\) to \(\operatorname{Hom}_{\mathbf{A}_{F}}(\mathbf{M}_{F}, \mathbf{N}_{F})\) be bijective, it is necessary and sufficient that \(j_{N}\) be bijective.

\begin{enumerate}
\def\labelenumi{(\alph{enumi})}
\setcounter{enumi}{5}
\item
  Take A to be the polynomial ring K{[}X, Y{]} over a commutative field K and F to be the (topologizing and idempotent) set of ideals of A containing a power of the ideal m = (X) + (Y). If M = A/AX and u: A → M is the canonical mapping, show that \(u_{F}: A_{F} \to M_{F}\) is not surjective. (Show that \(A_{F}\) is identified with A and \(M_{F}\) with K{[} \(\overline{Y}\) , 1/ \(\overline{Y}\) {]}, where \(\overline{Y}\) is the class of Y in M, so that X \(\overline{Y}\) = 0 and X(1/ \(\overline{Y}\) ) = 0.)
\item
  If we set \(\mathcal{F}^1\mathrm{M} = \mathrm{M}_{\mathcal{F}} / j_{\mathrm{M}}(\mathrm{M})\) , show that, for every exact sequence \(0\to \mathrm{M}\to \mathrm{N}\to \mathrm{P}\to 0\) of A-modules, there is an exact sequence
\end{enumerate}

\[
0 \rightarrow \mathcal {F} M \rightarrow \mathcal {F} N \rightarrow \mathcal {F} P \rightarrow \mathcal {F} ^ {1} M \rightarrow \mathcal {F} ^ {1} N \rightarrow \mathcal {F} ^ {1} P
\]

(use the snake diagram).

\begin{enumerate}
\def\labelenumi{(\alph{enumi})}
\setcounter{enumi}{7}
\tightlist
\item
  Let \(\mathcal{F}'\) be the set of left ideals \(l'\) of \(\mathbf{A}_{\mathcal{F}}\) such that \((\mathbf{A}_{\mathcal{F}})_s / l'\) is an \(\mathcal{F}\) -negligible A-module: show that \(\mathcal{F}'\) is the set of left ideals of \(\mathbf{A}_{\mathcal{F}}\) containing the \(j_{\mathbf{A}}(\mathfrak{m})\) , where \(\mathfrak{m} \in \mathcal{F}\) (observe that \(\mathfrak{m}_{\mathcal{F}} = \mathbf{A}_{\mathcal{F}}\) for all \(\mathfrak{m} \in \mathcal{F}\) , using the exact sequence of (d); consequently \(A_{\mathcal{F}}/j_{A}(m)\) is F-negligible). For every \(A_{F}\) -module \(M'\) , \(FM' = F'M'\) ( \(M'\) being considered as an A-module by means of \(j_{A}\) ), \(F'\) is topologizing and idempotent and \(M'_{F'}\) is canonically identified with \(M'_{F}\) .
\end{enumerate}

\begin{enumerate}
\def\labelenumi{\arabic{enumi}.}
\setcounter{enumi}{19}
\tightlist
\item
  Let \(\mathcal{F}\) be a topologizing and idempotent set of left ideals of a ring A.
\end{enumerate}

\begin{enumerate}
\def\labelenumi{(\alph{enumi})}
\item
  Suppose that every ideal of \(\mathcal{F}\) contains a finitely generated ideal belonging to \(\mathcal{F}\) . Let M be an A-module and \((\mathbf{N}_{\iota})_{\iota \in I}\) a right directed family of submodules of M, whose union is M. Show that \(\mathbf{M}_{\mathcal{F}}\) is the union of its sub- \(\mathbf{A}_{\mathcal{F}}\) -modules \((\mathbf{N}_{\iota})_{\mathcal{F}}\) . In particular, for every family \((\mathbf{M}_{\lambda})_{\lambda \in L}\) of A-modules, \(\left( \bigoplus_{\lambda \in L} \mathbf{M}_{\lambda} \right)_{\mathcal{F}}\) is canonically isomorphic to \(\bigoplus_{\lambda \in L} (\mathbf{M}_{\lambda})_{\mathcal{F}}\) .
\item
  Suppose that condition (a) holds and also that, for every surjective A-module homomorphism \(u \colon \mathbf{M} \to \mathbf{N}\) , \(u_{\mathcal{F}} \colon \mathbf{M}_{\mathcal{F}} \to \mathbf{N}_{\mathcal{F}}\) is surjective. Then, the canonical mapping \(\mathbf{A}_{\mathcal{F}} \otimes_{\mathbf{A}} \mathbf{M} \to \mathbf{M}_{\mathcal{F}}\) defined by the external law on \(\mathbf{M}_{\mathcal{F}}\) as an \(\mathbf{A}_{\mathcal{F}}\) -module is an \(\mathbf{A}_{\mathcal{F}}\) -module isomorphism (reduce it to the case where \(\mathbf{M}\) is free); the ring \(\mathbf{A}_{\mathcal{F}}\) is a flat left A-module and the left \(\mathbf{A}_{\mathcal{F}}\) -modules coincide with the left A-modules \(\mathbf{M}\) such that the mapping \(j_{\mathbf{M}} \colon \mathbf{M} \to \mathbf{M}_{\mathcal{F}}\) is bijective.
\item
  Suppose that A is an infinite product \(\biguplus_{i\in I}A_{i}\) of rings; let \(e_{i}\) denote the unit element of \(A_{i}\) and take F to be the set of left ideals of A containing the two-sided ideal \(a=\bigoplus_{i\in I}A_{i}\) (the \(A_{i}\) being canonically identified with ideals of A). Show that, for every left A-module M, \(M_{F}\) is identified with the product \(\prod_{i\in I}(e_{i}M)\) and in particular \(A_{F}=A\) . Deduce that, for every surjective A-homomorphism \(u: M\to N\) , \(u_{F}: M_{F}\to N_{F}\) is surjective, but give examples of A-modules M such that \(j_{M}\) is not bijective.
\item
  Suppose that condition (a) holds and also that A is commutative. Show that, if M is F-negligible, \(\mathbf{M}_{(\mathcal{F})}=0\) .
\end{enumerate}

\begin{enumerate}
\def\labelenumi{\arabic{enumi}.}
\setcounter{enumi}{20}
\tightlist
\item
  Let A be a commutative ring and \(\mathcal{F}\) a topologizing and idempotent set of ideals of A.
\end{enumerate}

\begin{enumerate}
\def\labelenumi{(\alph{enumi})}
\tightlist
\item
  Then the ring \(\mathbf{A}_{\mathcal{F}}\) is commutative (note that, if \(\mathfrak{m} \in \mathcal{F}\) and
\end{enumerate}

\[
u \in \operatorname{Hom} _ {\mathbf {A}} (\mathfrak {m}, \mathbf {A}),
\]

then \(u(\mathfrak{m}^2) \subset \mathfrak{m}\) ; if \(v\) is another element of \(\operatorname{Hom}_{\mathbf{A}}(\mathfrak{m}, \mathbf{A})\) , show that

\[
v (u (x y)) = u (v (x y))
\]

for \(x \in \mathfrak{m}, y \in \mathfrak{m}\) ).

\begin{enumerate}
\def\labelenumi{(\alph{enumi})}
\setcounter{enumi}{1}
\item
  Let \(\mathcal{F}'\) be the family of ideals \(l'\) of \(\mathbf{A}_{\mathcal{F}}\) such that \(\mathbf{A}_{\mathcal{F}} / l'\) is \(\mathcal{F}\) -negligible (exercise 19 (h)). Show that the mapping \(\mathfrak{p} \mapsto \mathfrak{p}_{\mathcal{F}}\) is a bijection of the set of ideals \(\mathfrak{p}\) of \(\mathbf{A}\) not belonging to \(\mathcal{F}\) onto the set of prime ideals of \(\mathbf{A}_{\mathcal{F}}\) not belonging to \(\mathcal{F}'\) . (Show first that, if \(\mathbf{A}\) is an integral domain, so is \(\Lambda_{\mathcal{F}}\) , making the same note as in (a); then use the exact sequence of Exercise 19 (d) to show that \(\mathfrak{p}_{\mathcal{F}}\) is prime. If \(\mathfrak{p}'\) is a prime ideal of \(A_{\mathcal{F}}\) not belonging to \(\mathcal{F}'\) and \(\mathfrak{p} = j_{\mathrm{A}}^{-1}(\mathfrak{p}')\) , recall that \(\mathfrak{p}' = \mathfrak{p}_{\mathcal{F}}\) .
\item
  Let B be an A-algebra and G the set of left ideals l of B such that \(B_{s}/!\) is F-negligible. Show that G is the set of left ideals of B containing an ideal of the form B.m, where \(m \in F\) ; deduce that G is topologizing and idempotent and that, for every B-module N, FN = GN and \(N_{F}\) is canonically identified with \(N_{G}\) .
\item
  Let S be a multiplicative subset of A and F the set of ideals of A meeting S. Show that F is topologizing and idempotent and that, for every A-module M, \(M_{(F)}\) and \(M_{F}\) are canonically identified with \(S^{-1}M\) (note that every ideal of F contains a principal ideal As, where \(s \in S\) ).
\end{enumerate}

¶ 22. Let A be a ring (not necessarily commutative) and S a multiplicative subset of A.

\begin{enumerate}
\def\labelenumi{(\alph{enumi})}
\item
  Let \(\mathcal{F}\) be the set of left ideals \(\mathfrak{l}\) of \(\mathbf{A}\) such that, for all \(a\in \mathbf{A}\) , there exists \(s\in \mathbf{S}\) such that \(sa\in \mathfrak{l}\) ; this implies in particular that \(\mathfrak{l}\cap \mathbf{S}\neq \varnothing\) . Show that \(\mathcal{F}\) is topologizing and idempotent.
\item
  Let B be a ring and \(\phi: A \to B\) a ring homomorphism. The ordered pair \((B, \phi)\) is called a ring of left fractions of \(A\) with denominators in \(S\) if it satisfies the following conditions:
\end{enumerate}

\begin{enumerate}
\def\labelenumi{(\Roman{enumi})}
\item
  If \(\phi(a) = 0\) , there exists \(s \in S\) such that \(sa = 0\) .
\item
  If \(s \in S\) , \(\phi(s)\) is invertible in \(B\) .
\item
  Every element of B is of the form \((\phi(s))^{-1}\phi(a)\) , where \(a \in \mathbf{A}\) .
\end{enumerate}

Show that the following properties are equivalent:

(α) The ring A possesses a ring of left fractions with denominators in S.

(β) The following conditions hold:

\((\beta_{1})\) For all \(s \in S\) , \(a \in A\) , there exists \(t \in S\) and \(b \in A\) such that \(ta = bs\) .

\((\beta_{2})\) If \(a \in \mathbf{A}\) and \(s \in \mathbf{S}\) satisfy \(as = 0\) , there exists \(t \in \mathbf{S}\) such that \(ta = 0\) .

(γ) The canonical images of the elements of S in \(A_{F}\) are invertible.

Moreover these properties imply the following:

(δ) The principal ideals As, where \(s \in S\) , belong to F and every ideal \(m \in F\) contains one of these principal ideals.

(ε) The left annihilator of every \(s \in S\) is contained in \(\mathcal{FA}_{s}\) .

(ζ) For every exact sequence \(0 \to \mathbf{M} \xrightarrow{h} \mathbf{N} \xrightarrow{p} \mathbf{P} \to 0\) of A-modules, the sequence \(0 \longrightarrow \mathbf{M}_{\mathcal{F}} \xrightarrow{h_{\mathcal{F}}} \mathbf{N}_{\mathcal{F}} \xrightarrow{p_{\mathcal{F}}} \mathbf{P}_{\mathcal{F}} \longrightarrow 0\) is exact.

(Show that \((\alpha) \Rightarrow (\beta) \Rightarrow (\gamma) \Rightarrow (\alpha)\) . To see that \((\gamma)\) implies \((\alpha)\) , note that, for every A-module M and every \(x \in \mathcal{FM}\) , there exists \(s \in S\) such that \(sx = 0\) . To see that \((\beta)\) implies \((\gamma)\) , prove first that \((\beta)\) implies \((\delta)\) , \((\varepsilon)\) and \((\zeta)\) . To show that \((\beta)\) implies \((\zeta)\) , establish, using \((\delta)\) , that, if \(m \in \mathcal{F}\) and

\[
u \in \operatorname{Hom} _ {\mathbf {A}} (\mathfrak {m}, \mathrm{P} / \mathscr {F P}),
\]

then there exists \(s \in S\) such that \(As \subset m\) and \(v \in \operatorname{Hom}_{\mathbf{A}}(As, N/\mathcal{F}N)\) such that the diagram

\[
\begin{array}{c} \text {As} \xrightarrow {v} \text {N / FN} \\ i \Big \downarrow \qquad \qquad \qquad \qquad \qquad \qquad \qquad \qquad \Big \downarrow p ^ {\prime} \\ \mathfrak {m} \xrightarrow [ u ]{} \text {P / FP} \end{array}
\]

is commutative, i being the canonical injection and \(p'\) the mapping derived from p by taking quotients. Consider finally, for all \(s \in S\) , the exact sequence

\[
0 \longrightarrow \mathrm{N} \longrightarrow \mathrm{A} \xrightarrow {\mu_ {s}} \mathrm{A} \longrightarrow \mathrm{A} / \mathrm{As} \longrightarrow 0
\]

where \(\mu_s\) is the homothety \(\lambda \mapsto \lambda s\) ; use \((\varepsilon)\) , \((\zeta)\) and Exercise 4 of Algebra, Chapter I, § 8.)

\begin{enumerate}
\def\labelenumi{(\alph{enumi})}
\setcounter{enumi}{2}
\item
  Deduce from (b) that, if A possesses a ring of left fractions \((\mathbf{B}, \phi)\) with denominators in S, this ring has the following universal property: for every ring homomorphism \(f: A \to C\) such that \(f(s)\) is invertible in C for all \(s \in S\) , there exists a unique homomorphism \(g: B \to C\) such that \(f = g \circ \phi\) . In particular B is canonically isomorphic to \(A_{F}\) .
\item
  The notion of a ring of right fractions of A with denominators in S is similarly defined. Deduce from (c) that, if there exist both a ring of left fractions and a ring of right fractions of A, these two rings are isomorphic.
\item
  Let E be a vector space over a field K, \((e_{t})_{t\in I}\) a basis of E and T the tensor algebra of E. Let J be a non-empty subset of I distinct from I and S the multiplicative subset of T generated by the \(e_{t}\) such that \(t\in J\) . Show that T admits no ring of left or right fractions with denominators in S.
\item
  With the notation of (e) suppose that I = N and consider the quotient ring A of T by the two-sided ideal generated by the elements \(e_{0}^{2}\) and \(e_{0}e_{i} - e_{i+1}e_{0}\) , for all i \textgreater{} 0, and \(e_{i}e_{j} - e_{j}e_{i}\) , for i \textgreater{} 0, j \textgreater{} 0. Let \(S'\) be the multiplicative set generated by the canonical images of the \(e_{i}\) in A, for i \textgreater{} 0; show that A admits a ring of left fractions, but not a ring of right fractions, with denominators in \(S'\) .
\item
  Suppose that A admits a ring of left fractions with denominators in S. In the set \(S \times A\) , let R be the relation between \((s, x)\) and \((t, y)\) : ``there exists \(r \in S\) , \(u \in A\) , \(v \in A\) such that r = us = vt and ux = vy''. Show that R is an equivalence relation and define on the set \(B = (S \times A)/R\) a ring structure such that the ordered pair consisting of S and the mapping \(\phi: A \to B\) , which maps \(x \in A\) to the class of \((1, x)\) , is a ring of left fractions of A with denominators in S (use the fact that, for \(s \in S\) and \(t \in S\) , there exists \(u \in A\) and \(v \in A\) such that r = us = tv belongs to S). Consider the case \(S = A - \{0\}\) (cf.~Algebra, Chapter I, §9, Exercise 8).
\end{enumerate}

\begin{enumerate}
\def\labelenumi{\arabic{enumi}.}
\setcounter{enumi}{22}
\tightlist
\item
  \begin{enumerate}
  \def\labelenumii{(\alph{enumii})}
  \tightlist
  \item
    Let A be a left Noetherian ring with no divisor of 0. Show that A possesses a field of left fractions (use Exercise 14 (a) to show that condition \((\beta_{1})\) of Exercise 22 (b) is fulfilled for \(S = A - \{0\}\) ).
  \end{enumerate}
\end{enumerate}

\begin{enumerate}
\def\labelenumi{(\alph{enumi})}
\setcounter{enumi}{1}
\tightlist
\item
  Let A be a left and right Noetherian ring with no divisor of 0. Show that every finitely generated left A-module M, each of whose elements \(\neq 0\) is free, is a submodule of a free A-module. (Embed M canonically in \(\mathbf{K} \otimes_{\mathbf{A}} \mathbf{M}\) , where K is the field of fractions (left or right) of A considered as a right A-module; if \((x_{i})\) is a basis of the left vector K-space \(\mathbf{K} \otimes_{\mathbf{A}} \mathbf{M}\) , show that there exists \(a \neq 0\) in A such that M is contained in the sub-A-module of \(\mathbf{K} \otimes_{\mathbf{A}} \mathbf{M}\) generated by the \(a^{-1}x_{i}\) .)
\end{enumerate}

¶ 24. Let A be a ring, commutative or not, and F the set of left ideals 1 of A such that \(l \cap m \neq 0\) for every left ideal \(m \neq 0\) of A.

In order that \(l \in F\) , it is necessary and sufficient that, for every element \(a \neq 0\) of A, there exist \(b \in A\) such that \(ba \neq 0\) and \(ba \in l\) . Every ideal l of F contains the left socle of A (Algebra, Chapter VIII, §5, Exercise 9) and, if A is Artinian, F consists of the ideals containing the left socle of A.

\begin{enumerate}
\def\labelenumi{(\alph{enumi})}
\tightlist
\item
  Show that \(\mathcal{F}\) is a topologizing set (Exercise 16). Deduce that the set \(\mathfrak{a}\) of \(a \in \mathbf{A}\) such that \(\mathfrak{l}a = 0\) for some \(\mathfrak{l} \in \mathcal{F}\) (union of the right annihilators of the \(\mathfrak{l} \in \mathcal{F}\) ) is a two-sided ideal of \(\mathbf{A}\) , which contains no idempotent. If \(\mathbf{A}\) is left Noetherian, show that \(\mathfrak{a}\) is nilpotent (prove first that every element \(a \in \mathfrak{a}\) is nilpotent, by considering the left annihilators of the \(a^n\) ; then use Exercise 26 (b) of Algebra, Chapter VIII, § 6).
\end{enumerate}

Let K be a commutative field and B the ring of polynomials with coefficients in K in indeterminates Y and \(X_{i}\) ( \(i \geqslant 1\) ). Let b be the ideal of B generated by the \(X_{i}X_{j}\) ( \(i \neq j\) ) and the \(Y^{i-1}X_{i}\) and let A be the quotient ring B/b; show that in A the ideal a defined above contains the class of Y, which is not nilpotent.

\begin{enumerate}
\def\labelenumi{(\alph{enumi})}
\setcounter{enumi}{1}
\item
  Show that, for every left ideal m of A, there exists a left ideal n of A such that \(m \cap n = 0\) and \(m + n \in F\) . (Take n to be an ideal which is maximal amongst the left ideals l such that \(m \cap l = 0\) .)
\item
  A ring A is called (left) neat if the two-sided ideal a defined in (a) is reduced to 0. If this is so, the set F is idempotent (let l, m, n be left ideals of A such that \(n \in F\) , \(l \subset n\) , \(n/l\) being F-negligible, and \(m \neq 0\) . If \(x \neq 0\) belongs to \(m \cap n\) , there exists \(r \in F\) such that \(r.x \subset l \cap m\) ).
\end{enumerate}

¶ 25. With the notation of Exercise 24, suppose that A is a left neat ring.

\begin{enumerate}
\def\labelenumi{(\alph{enumi})}
\item
  Show that the canonical mapping \(j_{\mathbf{A}}: \mathbf{A} \to \mathbf{A}_{\mathcal{F}}\) is injective, which allows us to identify A with a subgroup of \(\mathbf{A}_{\mathcal{F}}\) . Let \(\mathcal{F}'\) be the set of left ideals \(l'\) of \(\mathbf{A}_{\mathcal{F}}\) such that \(l'_{\mathcal{F}} = \mathbf{A}_{\mathcal{F}}\) ; show that \(\mathcal{F}'\) is the set of left ideals \(l'\) of \(\mathbf{A}_{\mathcal{F}}\) such that \(l' \cap m' \neq 0\) for every left ideal \(m' \neq 0\) of \(\mathbf{A}_{\mathcal{F}}\) . (Show first that, for every left ideal \(m' \neq 0\) of \(\mathbf{A}_{\mathcal{F}}\) , \(A \cap m' \neq 0\) ; note on the other hand that the ideals \(l' \in \mathcal{F}'\) are precisely the left ideals of \(\mathbf{A}_{\mathcal{F}}\) containing an ideal \(m \in \mathcal{F}\) .) Show that the ring \(\mathbf{A}_{\mathcal{F}}\) is left neat.
\item
  Show that \(\mathbf{A}_{\mathcal{F}}\) is an injective A-module, that is (Algebra, Chapter II, §2, Exercise 11) that, for every left ideal I of A and all
\end{enumerate}

\[
u \in \operatorname{Hom} _ {\mathbf {A}} (\mathfrak {l}, (\mathrm{A} _ {\mathcal {F}}) _ {s}),
\]

there exists \(a' \in \mathbf{A}_{\mathcal{F}}\) such that \(u(x) = xa'\) for all \(x \in \mathfrak{l}\) . (Use Exercise 24 (b) and Exercise 19 (c).) Deduce that every endomorphism of the A-module \(\mathbf{A}_{\mathcal{F}}\) is of form \(u: x \mapsto xa'\) where \(a' \in \mathbf{A}_{\mathcal{F}}\) (note that \(\operatorname{Hom}_{\mathbf{A}}(\mathbf{A}, \mathbf{A}_{\mathcal{F}}) = \operatorname{Hom}_{\mathbf{A}_{\mathcal{F}}}(\mathbf{A}_{\mathcal{F}}, \mathbf{A}_{\mathcal{F}})\) by Exercise 19 (e)); also \(\operatorname{Ker}(u) = (\operatorname{Ker}(u))_{\mathcal{F}}\) .

\begin{enumerate}
\def\labelenumi{(\alph{enumi})}
\setcounter{enumi}{2}
\item
  Show that, for every left ideal \(\mathfrak{l}\) of \(\mathbf{A}\) , \(\mathfrak{l}_{\mathcal{F}}\) is a direct factor of \(\mathbf{A}_{\mathcal{F}}\) . (Let \(\mathfrak{m}\) be a left ideal of \(\mathbf{A}\) such that \(\mathfrak{m} \cap \mathfrak{l} = 0\) and \(\mathfrak{m} + \mathfrak{l} \in \mathcal{F}\) ; extend the projection of \(\mathfrak{m} + \mathfrak{l}\) onto \(\mathfrak{m}\) , which is zero on \(\mathfrak{l}\) , to an endomorphism \(p\) of \(\mathbf{A}_{\mathcal{F}}\) ; show that \(p^2 = p\) and that \(\operatorname{Ker}(p) = \mathfrak{l}_{\mathcal{F}}\) .)
\item
  Show that \(A_{F}\) is an absolutely flat ring (Chapter I, §2, Exercise 17). (Reduce it to the case where \(A = A_{F}\) . If u is the endomorphism \(x \mapsto xa\) of the left A-module \(A_{s}\) , then \(\operatorname{Ker}(u) = (\operatorname{Ker}(u))_{\mathcal{F}}\) , and \(\operatorname{Ker}(u)\) therefore admits a complementary ideal m in \(A_{s}\) ; note that \(\operatorname{Im}(u)\) is isomorphic to m, hence injective (Algebra, Chapter II, §2, Exercise 11) and therefore a direct factor of \(A_{s}\) (loc. cit.). Finally, use Exercise 13 of Algebra, Chapter II, §1.)
\end{enumerate}

¶ 26. (a) Show that every Zorn ring A (Algebra, Chapter VIII, § 6, Exercise 13) with no nil ideal ≠0 is left neat; in particular, every absolutely flat ring and every primitive ring whose socle is not reduced to 0 is a left neat ring.

\begin{enumerate}
\def\labelenumi{(\alph{enumi})}
\setcounter{enumi}{1}
\item
  Let A be a left primitive ring whose socle is not reduced to 0; show that, if A is represented as a dense subring of a ring of endomorphisms \(\operatorname{End}_{\mathbb{D}}(\mathbf{V})\) of a vector space (Algebra, Chapter VIII, § 5, Exercise 10), the ring \(A_{F}\) corresponding to A (Exercise 25) is isomorphic to \(\operatorname{End}_{\mathbb{D}}(\mathbf{V})\) .
\item
  Show that every quasi-simple ring A (Algebra, Chapter VIII, § 5, Exercise 5) is left and right neat and that the corresponding ring \(\mathbf{A}_{\mathcal{F}}\) is quasi-simple. Deduce that, for every ring B not reduced to 0, there exists a homomorphism \(\phi: \mathrm{B} \to \mathrm{R}\) to a ring R which is quasi-simple, absolutely flat and such that \(\mathbb{R}_s\) is an injective R-module.
\item
  Let A be a left Noetherian ring with no nilpotent ideals. Show that A is left neat and that the corresponding ring \(A_{F}\) is a semi-simple ring (``Goldie's Theorem''; use Exercise 24 (a) and, on the other hand, use Algebra, Chapter VIII, § 2, Exercise 7, noting that in the absolutely flat ring \(A_{F}\) there cannot exist any family of idempotents which are pairwise orthogonal, unless A contained an infinite direct sum of left ideals \(\neq0\) ). Show that, if \(s\in A\) is not a right divisor of 0 in A, it is invertible in \(A_{F}\) and conversely; the set S of these elements is a multiplicative subset of A. Show that F is the set of left ideals of A which meet S (to see that the ideals of A meeting S belong to F, note that the elements of S are not left divisors of 0 in A; to see that every \(l\in F\) contains an element of S, represent each simple component of \(A_{F}\) as an endomorphism ring of an n-dimensional vector space E and define by induction n elements \(u_{1},\ldots,u_{n}\) such that \(u_{i}u_{j}=0\) for j\textless i and \(u_{i}^{2}\neq0\) , the \(u_{i}\) being of rank 1; finally show that the kernel of \(s=u_{1}+\cdots+u_{n}\) is zero). Deduce that \((\mathbf{A}_{\mathcal{F}},j_{\mathrm{A}})\) is a ring of left fractions of A with denominators in S (Exercise 22).
\end{enumerate}

\begin{enumerate}
\def\labelenumi{\arabic{enumi}.}
\setcounter{enumi}{26}
\tightlist
\item
  \begin{enumerate}
  \def\labelenumii{(\alph{enumii})}
  \tightlist
  \item
    Let A be a ring and M a finitely generated faithful A-module. Show that, if a, b are two ideals of A such that aM = bM, the radicals of a and b are equal (use no. 2, Corollary 2 to Proposition 4).
  \end{enumerate}
\end{enumerate}

\begin{enumerate}
\def\labelenumi{(\alph{enumi})}
\setcounter{enumi}{1}
\item
  Deduce from (a) that, if in an integral domain A p is a finitely generated prime ideal, then \(p^{m} \neq p^{n}\) for \(m \neq n\) .
\item
  Let K be a field and A the polynomial ring K{[}X, Y{]} in two indeterminates. Let a be the ideal \(AX^{4} + AX^{3}Y + AXY^{3} + AY^{4}\) , b the ideal \(a + AX^{2}Y^{2}\) of A; show that \(b \neq a\) and \(b^{2} = a^{2}\) .
\end{enumerate}
% semantic-fragments: legacy-input-seam
\relax{}
